\chapter{Justificación}
Realizar la justificación de la investigación para una tesis debe realizarse desde 5 perspectivas:

\begin{itemize}
    \item \textbf{Justificación Teórica:} Validar, utilizar, cuestionar (Valida teoría)
    \item \textbf{Justificación Práctica:} Hablar de (Resuelve el problema)
    \item \textbf{Justificación Metodológica:} Recrear (Utilizar una metdología existente, llevarlo a otro lugar y ajustarlo u adaptarlo a las necesidades de ese nuevo entorno) Replicar (Copiar y pegar el modelo a otra zona sin considerar el contexto). Presenta a la comunidad académica un modelo aplicativo que es una forma de resolver problemas de tales u cuáles naturalezas, el mismo que puede ser recreado en otras investigaciones.
    \item \textbf{Justificación Social:} Tiene que estar enmarcado dentro de las ODS, Responsabilidad Social Univesitaria y los Ejes estrátegicos del plan nacional al 2050. (Genera una impacto en la organización u la empresa)
    \item \textbf{Justificación Económica:} Impacto de nuestra investigación sobre la competitividad, economía de la organización. No se necesario que sea numérico; En términos generales debemos explicar que a través de nuestra tesis se logrará la estabilidad de la organización, la competititvidad de la empresa, se va a asegurar las sostenibilidad en el etiempo de la empresa en el sector u negocio en el cual estamos trabajando.
\end{itemize}

Unas 6 líneas por cada justificación es lo ideal. Las siguientes secciones incluyen un inicio de la justificación, esto es un ejemplo de cómo realicé la justificación de mi trabajo, pero cada trabajo de tesis es especial y único.

\section{Justificación Teórica}
    \begin{sectext}
        La presente investigación se consolida al establecer la base del mismo en un conjunto de teorías en distintos ámbitos de conocimiento;...
    \end{sectext}

\section{Justificación Metodológica}
    \begin{sectext}
        La presente investigación se enmarca al uso de 2 métodos consolidados en el ámbito de la creación e invención tecnológica en materia de hardware y software...
    \end{sectext}

\section{Justificación Práctica}
    \begin{sectext}
        La presente investigación se centra en la superación de la carente...
    \end{sectext}

\section{Justificación Social}
    \begin{sectext}
        La presente investigación enmarca su desarrollo sobre los Objetivos de Desarrollo Sostenible 3 (Salud y Bienestar) y 9 (Industria, innovación e infraestructuras) debido a sus implicancias en el ámbito médico por medio del ámbito ingenieril. Por su parte, respecto a la Responsabilidad Social Universitaria, la investigación abarca el área de Intervención Tecnológica debido a la disposición de la investigación a proponer y ejecutar tecnologías en favor de problemas sociales de la comunidad. A su vez, respecto al Plan Estatrégico de Desarrollo Nacional al 2050, la investigación toma especial énfasis en el ON3. Competitividad e Innovación, siendo el OE 3.5 Elevar la capacidad científica y de innovación tecnológica del país el objetivo al cual la investigación se alinea con mayor relevancia.
    \end{sectext}

\section{Justificación Económica}
    \begin{sectext}
        La presente investigación busca mejorar la capacidad de análisis del...
    \end{sectext}